\documentclass[10pt,a4paper]{amsart}
\usepackage[margin=19mm]{geometry}
\usepackage{amsmath,amssymb,mathtools}
\usepackage[scr=rsfs,cal=euler]{mathalfa}
\usepackage{xfrac}
\usepackage[unicode,hidelinks,bookmarks=false]{hyperref}
\newcommand{\ie}{i.e.\ }
\newcommand{\cf}{cf.\ }
\newcommand{\resp}{respectively\ }
\newcommand{\Subsection}{\subsection}
\providecommand{\nobreakdash}{\nobreak\hbox{-}}
\setlength{\emergencystretch}{2em}
\multlinegap0pt
\usepackage{amssymb}
\usepackage{tikz}
\usepackage{algorithm}\usepackage{algpseudocode}
\newtheorem{theorem}{Theorem}
\newtheorem{lemma}{Lemma}
\newcommand{\Ah}{A_h}
\title{Numerical identification of the time-dependent coefficient in the heat equation with fractional Laplacian}
\author{Arshyn Altybay, Gulzat Nalzhupbayeva}
\date{Source: arXiv:2511.16238v2 (CC BY 4.0)}
\begin{document}
\maketitle

\noindent\emph{Source and licence.} Numerical identification of the time-dependent coefficient in the heat equation with fractional Laplacian. Version arXiv:2511.16238v2 (CC BY 4.0), distributed by arXiv under the Creative Commons Attribution 4.0 International licence, http://creativecommons.org/licenses/by/4.0/, as recorded in the arXiv licence metadata for this exact version and shown on the abstract page https://arxiv.org/abs/2511.16238v2. Full licence notice: \texttt{licenses/arxiv-2511.16238-CC-BY-4.0.txt}.\par\smallskip

\noindent\textit{Editorial scope.} Counted unit: the article's theorem thm:inverse-convergence (source0.tex lines 1147--1213). The article's complete original proof of it is delivered with it (source0.tex lines 1215--1536, one \texttt{proof} environment of 322 lines). No mathematics is written, repaired or re-derived: the statement and the proof are the article's own lines, in the article's own notation, and no retained line is substituted or re-flowed.  Retained uncounted as the source-local context the proof needs: lines 541--545; lines 659--708; lines 751--760; lines 930--1006.  Nothing was omitted from any retained span, so the package carries no omission record.  No retained line contains a citation, so the wrapper carries no bibliography and no bibliography entry.  No retained line contains a figure environment, an image-inclusion command or drawing code, so no image is delivered and the package declares no asset dependency.  Editorial formatting only, in addition: an AMS \texttt{amsart} preamble that restates the article's own declarations and macros used by the retained lines, namely the environments \texttt{theorem}, \texttt{lemma} on separate counters, together with the standard packages those lines need.  The retained environments print in this document as Theorem 1, Lemma 1 in order of appearance, whereas the article prints the same items with the numbering of its own full document.  The complete native source of the article as distributed in the arXiv e-print for this exact version is bundled unchanged as \texttt{sources/189502/source0.tex}.
\begin{equation}\label{fds_full}
 \frac{\mathbf U^{n+1}-\mathbf U^n}{\tau}
 +\Ah\frac{\mathbf U^{n+1}+\mathbf U^n}{2}
 =r^{n+\frac12}\mathbf F^{n+\frac12},
\end{equation}

\begin{theorem}[Second-order convergence of the direct scheme]
\label{thm:direct-convergence}
Assume that the exact solution \(u\) is sufficiently regular and that
\begin{equation}
\sup_{0\le t\le T}
\left(
\|(u_{ttt}(t))_h\|_h
+
\|((rf)_{tt}(t))_h\|_h
\right)
\le C_t,
\label{eq:time-regularity-nodal}
\end{equation}
uniformly with respect to \(h\).
Assume further that the spatial consistency estimate
\begin{equation}
\|A_h v_h-(\mathcal L^s v)_h\|_h
\le C_v h^2
\label{eq:spatial-consistency-thm}
\end{equation}
holds uniformly for
\[
v=u(t,\cdot),\qquad 0\le t\le T,
\]
with a constant independent of \(h\).

Let \(U^n\) be the solution of the Crank--Nicolson scheme
\eqref{fds_full}
with exact initial data
\[
U^0=u_h^0.
\]
Then there exists a constant \(C>0\), independent of \(h\) and \(\tau\),
such that
\begin{equation}
\max_{0\le n\le M}
\|u_h^n-U^n\|_h
\le
C(\tau^2+h^2),
\label{eq:L2-direct-convergence}
\end{equation}
and
\begin{equation}
\max_{0\le n\le M}
\|u_h^n-U^n\|_{A_h,h}
\le
C(\tau^2+h^2).
\label{eq:energy-direct-convergence}
\end{equation}
\end{theorem}

\begin{equation}
\frac{u_h^{n+1}-u_h^n}{\tau}
+
A_h\frac{u_h^{n+1}+u_h^n}{2}
=
r^{n+\frac12}F^{n+\frac12}
+
\xi^{n+\frac12},
\label{eq:exact-CN-relation-revised}
\end{equation}

\begin{lemma}[Discrete identifiability]
\label{lem:discrete-identifiability}
Assume the continuous nondegeneracy condition
\begin{equation}
|(f(t),\omega)_{L^2}|
\ge
\kappa_0>0,
\qquad
0\le t\le T.
\label{eq:continuous-nondegeneracy-recall}
\end{equation}
Suppose further that
\begin{equation}
\left|
\langle F^{n+\frac12},\omega\rangle_h
-
(f(t_{n+\frac12}),\omega)_{L^2}
\right|
\le
C_Qh^2,
\label{eq:quadrature-fomega}
\end{equation}
uniformly in \(n\), and that
\begin{equation}
\sup_{n,h}\|F^{n+\frac12}\|_h\le C_F,
\qquad
\sup_h\|A_h\omega\|_h\le C_\omega.
\label{eq:identifiability-bounds}
\end{equation}

Let
\[
S^{n+\frac12}
=
\left(I+\frac{\tau}{2}A_h\right)^{-1}
F^{n+\frac12},
\]
and define
\[
d_h^{n+\frac12}
=
\langle S^{n+\frac12},\omega\rangle_h.
\]
Then
\begin{equation}
\left|
d_h^{n+\frac12}
-
(f(t_{n+\frac12}),\omega)_{L^2}
\right|
\le
C_Qh^2
+
\frac{\tau}{2}C_FC_\omega.
\label{eq:discrete-denominator-consistency}
\end{equation}
Consequently, if \(h\) and \(\tau\) are sufficiently small so that
\begin{equation}
C_Qh^2
+
\frac{\tau}{2}C_FC_\omega
\le
\frac{\kappa_0}{2},
\label{eq:mesh-identifiability-condition}
\end{equation}
then
\begin{equation}
|d_h^{n+\frac12}|
\ge
\frac{\kappa_0}{2}
=:\kappa
>0,
\qquad
n=0,\ldots,M-1.
\label{eq:discrete-identifiability}
\end{equation}
\end{lemma}

\begin{theorem}[Conditional convergence of the coupled reconstruction]
\label{thm:inverse-convergence}
Assume the hypotheses of
Theorem~\ref{thm:direct-convergence} and
Lemma~\ref{lem:discrete-identifiability}.
Assume in addition that
\begin{equation}
\sup_h\|A_h\omega\|_h\le C_\omega,
\qquad
\sup_{n,h}\|F^{n+\frac12}\|_h\le C_F,
\label{eq:inverse-uniform-bounds}
\end{equation}
and that
\[
w\in C^3[0,T].
\]

Suppose that the derivative approximation constructed from the
measurement data satisfies
\begin{equation}
\left|
z^{n+\frac12}
-
w'(t_{n+\frac12})
\right|
\le
\eta,
\qquad
n=0,\ldots,M-1.
\label{eq:derivative-error-revised}
\end{equation}
Assume also that the nodal quadrature of the differentiated
measurement satisfies
\begin{equation}
\left|
w'(t)
-
\left\langle
(u_t(t,x_i))_{i=1}^{N-1},
\omega
\right\rangle_h
\right|
\le
C_Qh^2,
\qquad
0\le t\le T.
\label{eq:measurement-quadrature-revised}
\end{equation}

Then, for sufficiently small \(h\) and \(\tau\), there exists a
constant \(C>0\), independent of \(h\), \(\tau\), and \(\eta\), such
that
\begin{equation}
\max_{0\le n\le M}
\|u_h^n-U^n\|_h
+
\max_{0\le n\le M-1}
\left|
r(t_{n+\frac12})
-
r_h^{n+\frac12}
\right|
\le
C(\tau^2+h^2+\eta).
\label{eq:coupled-convergence-revised}
\end{equation}
\end{theorem}

\begin{proof}
Set
\[
e^n=u_h^n-U^n,
\qquad
\rho^{n+\frac12}
=
r(t_{n+\frac12})-r_h^{n+\frac12}.
\]

For brevity, write
\[
r_*^{n+\frac12}:=r(t_{n+\frac12}),
\qquad
L=I+\frac{\tau}{2}A_h,
\qquad
R=I-\frac{\tau}{2}A_h.
\]
The exact nodal solution satisfies
\begin{equation}
Lu_h^{n+1}
=
Ru_h^n
+
\tau r_*^{n+\frac12}F^{n+\frac12}
+
\tau\xi^{n+\frac12},
\label{eq:exact-state-LR}
\end{equation}
where
\begin{equation}
\|\xi^{n+\frac12}\|_h
\le
C(\tau^2+h^2)
\label{eq:inverse-defect-bound}
\end{equation}
by Theorem~\ref{thm:direct-convergence}.

Define
\[
Y_*^n=L^{-1}Ru_h^n,
\qquad
S^{n+\frac12}=L^{-1}F^{n+\frac12},
\qquad
V_*^n=\frac{u_h^n+Y_*^n}{2}.
\]
Equation \eqref{eq:exact-state-LR} gives
\[
u_h^{n+1}
=
Y_*^n
+
\tau r_*^{n+\frac12}S^{n+\frac12}
+
\tau L^{-1}\xi^{n+\frac12}.
\]
Therefore
\begin{equation}
\frac{u_h^{n+1}+u_h^n}{2}
=
V_*^n
+
\frac{\tau}{2}
r_*^{n+\frac12}S^{n+\frac12}
+
\frac{\tau}{2}L^{-1}\xi^{n+\frac12}.
\label{eq:exact-midpoint-state}
\end{equation}

We next derive the corresponding scalar identity for the exact
coefficient. Taking the discrete inner product of
\eqref{eq:exact-CN-relation-revised} with \(\omega\) gives
\begin{equation}
\left\langle
\frac{u_h^{n+1}-u_h^n}{\tau},
\omega
\right\rangle_h
+
\left\langle
A_h\frac{u_h^{n+1}+u_h^n}{2},
\omega
\right\rangle_h
=
r_*^{n+\frac12}
\langle F^{n+\frac12},\omega\rangle_h
+
\langle\xi^{n+\frac12},\omega\rangle_h.
\label{eq:exact-measurement-discrete}
\end{equation}

Since
\[
\frac{1}{\tau}
\left\langle
u_h^{n+1}-u_h^n,\omega
\right\rangle_h
=
\frac{1}{\tau}
\int_{t_n}^{t_{n+1}}
\langle (u_t(t,x_i))_i,\omega\rangle_h\,dt,
\]
assumption \eqref{eq:measurement-quadrature-revised} implies
\[
\left|
\frac{1}{\tau}
\left\langle
u_h^{n+1}-u_h^n,\omega
\right\rangle_h
-
\frac{w(t_{n+1})-w(t_n)}{\tau}
\right|
\le
C h^2.
\]
Since \(w\in C^3[0,T]\),
\[
\frac{w(t_{n+1})-w(t_n)}{\tau}
=
w'(t_{n+\frac12})
+
O(\tau^2).
\]
Hence
\begin{equation}
\left|
\frac{1}{\tau}
\left\langle
u_h^{n+1}-u_h^n,\omega
\right\rangle_h
-
w'(t_{n+\frac12})
\right|
\le
C(\tau^2+h^2).
\label{eq:measurement-difference-bound}
\end{equation}

Using \eqref{eq:exact-midpoint-state},
the symmetry of \(A_h\), and
\[
F^{n+\frac12}
=
S^{n+\frac12}
+
\frac{\tau}{2}A_hS^{n+\frac12},
\]
equation \eqref{eq:exact-measurement-discrete} can be rearranged as
\begin{equation}
r_*^{n+\frac12}
d_h^{n+\frac12}
=
w'(t_{n+\frac12})
+
\langle A_hV_*^n,\omega\rangle_h
+
\zeta^{n+\frac12},
\label{eq:exact-coefficient-identity}
\end{equation}
where
\begin{equation}
|\zeta^{n+\frac12}|
\le
C(\tau^2+h^2).
\label{eq:zeta-bound}
\end{equation}
Indeed, the terms contributing to \(\zeta^{n+\frac12}\) are
the measurement approximation error
\eqref{eq:measurement-difference-bound},
the defect term
\(\langle\xi^{n+\frac12},\omega\rangle_h\),
and
\[
\frac{\tau}{2}
\langle
A_hL^{-1}\xi^{n+\frac12},
\omega
\rangle_h.
\]
For the last term, symmetry gives
\[
\left|
\left\langle
A_hL^{-1}\xi^{n+\frac12},
\omega
\right\rangle_h
\right|
=
\left|
\left\langle
L^{-1}\xi^{n+\frac12},
A_h\omega
\right\rangle_h
\right|
\le
C_\omega
\|\xi^{n+\frac12}\|_h,
\]
because \(\|L^{-1}\|_h\le1\).

The numerical coefficient satisfies
\begin{equation}
r_h^{n+\frac12}
d_h^{n+\frac12}
=
z^{n+\frac12}
+
\langle A_hV^n,\omega\rangle_h,
\label{eq:numerical-coefficient-identity}
\end{equation}
where
\[
V^n
=
\frac{U^n+L^{-1}RU^n}{2}.
\]
Subtracting
\eqref{eq:numerical-coefficient-identity}
from \eqref{eq:exact-coefficient-identity} gives
\[
\rho^{n+\frac12}d_h^{n+\frac12}
=
w'(t_{n+\frac12})-z^{n+\frac12}
+
\langle A_h(V_*^n-V^n),\omega\rangle_h
+
\zeta^{n+\frac12}.
\]
By symmetry,
\[
\left|
\langle A_h(V_*^n-V^n),\omega\rangle_h
\right|
\le
\|V_*^n-V^n\|_h
\|A_h\omega\|_h.
\]
Moreover,
\[
V_*^n-V^n
=
\frac12
\left(
I+L^{-1}R
\right)e^n.
\]
Since
\[
\|L^{-1}R\|_h\le1,
\]
we obtain
\[
\|V_*^n-V^n\|_h
\le
\|e^n\|_h.
\]
Lemma~\ref{lem:discrete-identifiability} gives
\[
|d_h^{n+\frac12}|
\ge\kappa>0.
\]
Consequently,
\begin{equation}
|\rho^{n+\frac12}|
\le
C\left(
\|e^n\|_h
+
\tau^2+h^2+\eta
\right).
\label{eq:rho-estimate-revised}
\end{equation}

It remains to estimate the state error. Subtracting the numerical
state equation from \eqref{eq:exact-state-LR} yields
\[
Le^{n+1}
=
Re^n
+
\tau\rho^{n+\frac12}F^{n+\frac12}
+
\tau\xi^{n+\frac12}.
\]
Since
\[
\|L^{-1}R\|_h\le1,
\qquad
\|L^{-1}\|_h\le1,
\]
and \(\|F^{n+\frac12}\|_h\le C_F\),
\[
\|e^{n+1}\|_h
\le
\|e^n\|_h
+
C\tau|\rho^{n+\frac12}|
+
C\tau(\tau^2+h^2).
\]
Using \eqref{eq:rho-estimate-revised},
\[
\|e^{n+1}\|_h
\le
(1+C\tau)\|e^n\|_h
+
C\tau(\tau^2+h^2+\eta).
\]
Since \(e^0=0\), the discrete Gr\"onwall inequality yields
\[
\max_{0\le n\le M}\|e^n\|_h
\le
C(\tau^2+h^2+\eta).
\]
Substitution into \eqref{eq:rho-estimate-revised} gives
\[
\max_{0\le n\le M-1}
|\rho^{n+\frac12}|
\le
C(\tau^2+h^2+\eta).
\]
This proves \eqref{eq:coupled-convergence-revised}.
\end{proof}

\end{document}
